% Hierarchical MARL for relay-drone delivery in disaster response (English version)
% Target version: v3 (26_08_31_19)
% Compile: pdflatex
\documentclass[12pt]{article}
\usepackage[margin=1in]{geometry}
\usepackage{setspace}
\usepackage{amsmath, amssymb}
\usepackage{booktabs}
\usepackage{graphicx}
\usepackage{url}
\doublespacing
\pagestyle{empty}

\title{Hierarchical Multi-Agent Reinforcement Learning for\\
Relay-Drone Network Delivery in Disaster Response}
\author{Kwangjong Ko \and Hyunseok Choi \and Taesu Cheong\thanks{Corresponding author.}\\
Department of Industrial and Management Engineering, Korea University}
\date{}

\begin{document}
\maketitle

\begin{abstract}
In disaster zones where communication infrastructure has been destroyed, the direct communication
range of a drone bounds the region it can serve. A relay drone network, in which each drone
forwards its peers' traffic, relaxes that bound but introduces an operational tension: a drone
holding a relay position cannot deliver. This study addresses the setting in which a homogeneous
fleet delivers relief supplies while simultaneously forming a multi-hop relay chain so that every
aircraft remains connected to the control center at all times. The objective is to minimize the
makespan of serving all demand points.

Treating connectivity as a reward penalty was found to inject unavoidable penalties into the
learning signal: a drone stationary for unloading can be stranded by teammates departing, and no
action available to it prevents the resulting penalty. Connectivity is therefore enforced as a
hard environmental constraint through a minimal-perturbation projection that shrinks each commanded
displacement to the largest value preserving connectivity. Separately, single-level continuous
control was observed not to acquire the reload cycle---returning to the depot once capacity is
exhausted---so a hierarchy is proposed in which a manager selects a destination, a return, or a
relay hold as a discrete action, and a worker produces two-dimensional thrust.

Evaluation used 60 held-out instances rolled out five times each, for 300 rollouts, at
communication range 300 with 50 destinations and 4 drones, and training was repeated under four
random seeds. The proposed method served 24.90 demand points (standard error across seeds 0.50,
95\% confidence interval $[23.32, 26.48]$) against 23.04 for a greedy heuristic, a difference that
is statistically significant ($t=3.74$, $p=0.034$, with all four seeds exceeding the baseline).
Differences in the secondary metrics are larger: maximum reach from the control center rose from
649 to 807--830, a 25\% increase reaching roughly 2.7 times the direct communication range, and
movements cancelled by the connectivity constraint fell by more than 60\%, from 3,341 to
1,051--1,434. Fine-tuning both levels jointly degraded performance to 21.43, which is attributed
to the manager facing a non-stationary environment as worker updates change what its chosen
subgoals mean. No connectivity loss occurred during any training episode or evaluation rollout.
\end{abstract}

\section{Introduction}

Ground transport is frequently unusable in disaster response because road infrastructure is
damaged, and drone delivery has accordingly been studied as an alternative. What the logistics
literature has largely left aside is that communication infrastructure is damaged as well.
Relief cargo carries a high completion priority, so its delivery status must be tracked, and
loading and unloading demand fine control supported by real-time video. Continuous connectivity
during the mission is therefore a requirement rather than a convenience.

A relay drone network, in which each aircraft forwards the traffic of others, extends the fleet's
operating radius beyond the direct communication range. This structure creates an operational
trade-off. A drone that holds a relay position performs no deliveries, so the number of aircraft
assigned to relaying is exchanged directly against delivery throughput. A preliminary sweep
reported in Section~5.1 shows a regime in which the optimal point of this exchange differs across
instances; that regime is where fixed rules cannot be optimal and adaptive assignment has room to
contribute.

Three contributions are offered. First, sweeping problem difficulty along the communication range
identifies both where rule-based policies begin to fail and where no fixed role split can be
optimal. Second, enforcing connectivity by minimal-perturbation projection rather than by reward
removes connectivity loss structurally. Third, a hierarchical multi-agent reinforcement learning
framework separating destination selection from low-level maneuvering is proposed; it exceeds a
greedy heuristic by 8.1\% in delivery volume while extending reach by 25\% and reducing constraint
collisions by more than 60\%. Joint training of both levels is shown to degrade performance,
establishing the need for staged learning.

\section{Related Work}

\subsection{Drone delivery routing}

Murray and Chu (2015) formulated the flying sidekick traveling salesman problem, in which one truck
and one drone cooperate, and proposed a mixed-integer program together with heuristics that
minimize makespan. Murray and Raj (2020) extended the formulation to a heterogeneous fleet of
drones with explicit launch and retrieval queueing. Agatz et al. (2018) permitted launch and
retrieval at the same node and contributed an exact dynamic-programming procedure. Dorling et al.
(2017) defined a drone-only multi-trip vehicle routing problem with an energy model linear in
payload and battery weight. This line of work assumes that communication is free and unmodeled.

\subsection{UAV relay networks}

Zeng et al. (2016) proposed mobile relaying, in which a UAV serves as a moving relay node, and
optimized power allocation and trajectory alternately, although a single relay UAV was assumed.
Yanmaz (2022) decoupled mission UAVs from relay UAVs and computed relay placements from Steiner
points so that a minimal number of relays maintains team connectivity. Ding et al. (2022) jointly
determined trajectory, spectrum allocation, and next-hop selection in a multi-hop UAV relay network
using QMIX-based multi-agent reinforcement learning. Objectives in this line are throughput and
coverage; payload capacity and reload trips are not represented.

\subsection{The gap}

Role switching between task execution and communication relaying is itself established in the
multi-robot exploration literature. Cesare et al. (2015) defined explore, meet, sacrifice, and
relay states, with relay robots landing and remaining stationary to sustain data flow. Role
switching therefore cannot be claimed as novel here. What was not located within the scope of this
survey is the transplantation of that mechanism into a \emph{capacitated, multi-trip,
makespan-minimizing} delivery setting, where the relay obligation competes directly with payload
throughput and the fleet must repeatedly return to a depot to reload.

\subsection{Learning methodology}

Huang and Ontañón (2022) compared suppressing infeasible actions by reward penalty against blocking
them by masking under controlled conditions, and showed that episodic return under the penalty
treatment collapses as problem size grows. Dalal et al. (2018) introduced a safety layer that
analytically computes the smallest action correction satisfying the constraints. Ng et al. (1999)
proved that potential-based shaping is the necessary and sufficient form preserving the optimal
policy, and gave cases in which a positive per-timestep reward not of that form induces cycling
rather than task completion. Zhou et al. (2019) showed that representing rotation by a scalar angle
is discontinuous and that the continuous representation of SO(2) is $[\cos\theta, \sin\theta]$.
The present design follows these results.

\section{Problem Statement}

One control center and $|C|$ demand points are placed on a $1000 \times 1000$ plane. A homogeneous
fleet of $|K|$ drones departs the control center carrying $Q$ units. When drone $i$ enters the
interaction radius $r$ of demand point $j$, a delivery occurs, capacity decreases by one, and $S$
steps of unloading time elapse. A drone with zero capacity that returns to the control center is
restored to $Q$ over the same $S$ steps.

The communication graph $G_t = (V, E_t)$ is defined by drone positions at time $t$; two nodes are
adjacent when separated by no more than the communication range $R$, and the control center is a
node. The constraint is
\begin{equation}
\forall t, \ \forall i \in K: \ i \text{ is connected to the control center in } G_t .
\end{equation}
The objective is to minimize $T^*$, the time at which all demand points have been served.

A remark on mathematical programming is warranted. In a node-based routing formulation a drone in
transit has no coordinates, so connectivity cannot be evaluated at every instant. Closing that gap
requires either discretizing space or pre-specifying candidate hovering points, and the latter is
itself a placement problem. Because the same drones both deliver and relay, placement cannot be
separated into a subproblem as in Yanmaz (2022), and the result is a coupled placement-and-routing
problem.

\section{Method}

\subsection{Connectivity as a hard constraint}

In the preceding implementation, connectivity was handled by assigning a large penalty on
disconnection and terminating the episode. A defect of this treatment was identified. When a drone
stationary for unloading is stranded by departing teammates, it receives a penalty that no action
available to it could have prevented. Unavoidable penalties corrupt the learning signal.

Connectivity is instead enforced as a hard environmental constraint. Given commanded displacements
$\Delta$, a scale vector $s \in [0,1]^{|K|}$ is selected so that the resulting configuration
$x + \Delta \odot s$ leaves every drone connected. Starting from $s = \mathbf{1}$, while the
configuration is disconnected the drone whose restriction yields the largest number of connected
drones is selected, and its scale is reduced by bisection to the largest safe value. Ties are
broken at random; breaking them deterministically was observed to concentrate the restriction
repeatedly on the drone anchoring the relay chain. Because reverting all scales to zero recovers
the previous configuration, which is connected by induction, the procedure always terminates in a
connected configuration.

\subsection{Hierarchy}

Single-level continuous control did not acquire the cycle of exhausting capacity, returning, and
reloading (Section~5.3). That cycle spans hundreds of steps. Decision making is therefore split
across two levels.

\textbf{Manager.} Every $H$ steps, and immediately whenever the current subgoal is invalidated,
one of $K+2$ options is chosen per drone: one of the $K$ nearest unserved destinations, a return to
the control center, or holding the current position. Admitting all destinations as actions would
tie the action space to instance size, so the candidate set is fixed at the $K$ nearest. Including
hold as an explicit action is central to the design: the relay role consists precisely of not
moving, and only an explicit option lets the manager allocate roles directly. The manager receives
only the team objective---deliveries, reloads, and a time penalty.

\textbf{Worker.} At every step a two-dimensional thrust $(v_x, v_y) \in [-1,1]^2$ is emitted,
normalized when its norm exceeds one, giving displacement $v \cdot v_{\max}$. A scalar-angle
representation is a discontinuous representation of SO(2), placing physically adjacent headings at
opposite ends of the action range, and reaching maximum speed requires tanh saturation, which the
maximum-entropy objective suppresses; the measured mean commanded speed of an angle-parameterized
policy was 0.511 of maximum. The worker reward comprises progress toward its subgoal, an arrival
bonus, and a penalty proportional to displacement removed by the connectivity correction. It
receives no team reward.

Both levels are trained with soft actor-critic under centralized training with decentralized
execution. The manager uses discrete SAC over a categorical policy; the worker uses continuous SAC
with a tanh-Gaussian policy.

\section{Experiments}

\subsection{Problem difficulty}

Greedy performance was measured while sweeping the communication range (Table~\ref{tab:sweep}).
The final column reports the gap between the best \emph{fixed} number of relay drones and an
oracle that selects the best number per instance.

\begin{table}[h]
\centering
\caption{Greedy performance by communication range (12 instances, deliveries out of 50)}
\label{tab:sweep}
\begin{tabular}{rrrrrrrr}
\toprule
Range & 0 relays & 1 & 2 & 3 & Best fixed & Per-instance oracle & Gap \\
\midrule
500 & 49.3 & 50.0 & 50.0 & 36.4 & 50.0 & 50.0 & $+0.0$ \\
400 & 40.6 & 48.4 & 43.9 & 17.9 & 48.4 & 49.4 & $+1.0$ \\
\textbf{300} & 21.0 & 22.7 & 18.5 & 9.3 & 22.7 & \textbf{28.2} & $\mathbf{+5.6}$ \\
250 & 14.4 & 15.2 & 12.6 & 7.8 & 15.2 & 19.1 & $+3.8$ \\
200 & 10.8 & 9.7 & 8.2 & 5.9 & 10.8 & 11.8 & $+1.0$ \\
150 & 7.0 & 5.8 & 4.6 & 4.2 & 7.0 & 7.2 & $+0.2$ \\
\bottomrule
\end{tabular}
\end{table}

Through range 400 the greedy policy essentially solves the problem, leaving no room for a learning
approach to contribute. At range 300 performance falls by more than half while the gap reaches its
maximum, indicating that no fixed relay count suffices and that the required allocation differs
across instances. Below range 200 the gap narrows again because no allocation performs well.
Subsequent experiments use range 300.

\subsection{Main results}

Thirty held-out instances were each evaluated three times, for 90 rollouts per method
(Table~\ref{tab:main}). Repetition averages over the randomness in the connectivity projection's
tie-breaking.

\begin{table}[h]
\centering
\caption{Main results (60 held-out instances $\times$ 5 repetitions = 300 rollouts, range 300)}
\label{tab:main}
\begin{tabular}{lrrrrrr}
\toprule
Method & Delivered & SD & SE & Reloads & Max reach & Moves cancelled \\
\midrule
Greedy & 23.04 & 8.89 & 0.51 & 2.11 & 649 & 3{,}341 \\
Greedy + 1 fixed relay & 14.51 & 6.11 & 0.35 & 1.06 & 533 & 2{,}635 \\
Proposed (seed 0) & 24.72 & 6.38 & 0.37 & 2.37 & 807 & 1{,}324 \\
Proposed (seed 1) & \textbf{25.95} & 6.61 & 0.38 & 2.92 & 826 & 1{,}308 \\
Proposed (seed 2) & 23.61 & 6.86 & 0.40 & 2.49 & 830 & 1{,}434 \\
Proposed (seed 3) & 25.32 & 5.68 & 0.33 & 2.50 & 808 & 1{,}051 \\
\textbf{Proposed (mean of 4 seeds)} & \textbf{24.90} & --- & \textbf{0.50} & 2.57 & \textbf{818} &
\textbf{1{,}279} \\
Joint training (both levels) & 21.43 & 5.20 & 0.30 & 2.30 & 680 & 1{,}300 \\
\bottomrule
\end{tabular}
\end{table}

Treating the four training seeds as the sample, the proposed method averages 24.90 deliveries with
a seed-level standard deviation of 1.00 and standard error of 0.50. The 95\% confidence interval
$[23.32,\,26.48]$ excludes the greedy value of 23.04, and a one-sample $t$ test gives $t=3.74$,
$p=0.034$. All four seeds exceeded the baseline, and the improvement is 8.1\%.

Differences in the secondary metrics are larger. Maximum reach from the control center rose from
649 to 807--830; with a communication range of 300, the relay chain extended to roughly 2.7 times
the direct range. Movements cancelled by the connectivity constraint fell by more than 60\%, from
3,341 to 1,051--1,434. The proposed method therefore travels farther while colliding less with the
constraint, and reloads rose to 2.37--2.92 against 2.11 for greedy.

Fine-tuning both levels jointly, by contrast, degraded performance to 21.43, below the greedy
baseline. The same degradation reproduced under a different seed, and both hop-depth ratio (0.50)
and reach (680) fell below the staged configuration. Because updating the worker changes what the
manager's chosen subgoals mean, the manager faces a non-stationary environment and its learned
assignment collapses. The joint stage is excluded from the final configuration.

No connectivity loss occurred in any training episode or evaluation rollout, confirming that the
hard constraint behaves as designed.

\subsection{Contribution of the hierarchy}

The bottleneck the hierarchy targeted was the reload cycle. Policies learned under a single-level
architecture performed between 0.0 and 0.4 reloads per episode. With 4 drones of capacity 5,
20 deliveries is a structural ceiling without reloading, and observed delivery volumes sat below
it. Under the hierarchy, reloads recovered to 2.34 per episode, comparable to greedy at 2.12.

Sensitivity to the manager decision period $H$ was examined, and $H=20$ dominated $H=10$.
Frequent re-decision replaces subgoals before the worker reaches them.

\subsection{Ablations}

Several design elements were tested individually. Potential-based shaping on relay spacing,
weighted redistribution of team reward toward drones necessary for relaying, a floor on the entropy
coefficient, and expanding observations to all peers each failed to improve delivery performance.
Removing the clustering penalty was preferable, and it was excluded from the final design.

Changing the action representation did not alter delivery performance but affected training
stability markedly. Held-out performance of the angle-parameterized policy declined 40\% after its
peak, whereas the vector-parameterized policy declined 8\%, and its peak occurred four times later
in training.

\section{Conclusion}

This study approached relay-drone delivery under disaster conditions with hierarchical multi-agent
reinforcement learning. Enforcing connectivity by minimal-perturbation projection removed
connectivity loss structurally, and separating destination selection from low-level maneuvering
recovered the reload cycle that single-level control had failed to acquire.

In delivery throughput the proposed method exceeded the greedy heuristic by 8.1\%, reproduced
across four training seeds ($p=0.034$). The margin itself is modest, but the difference in relay
network utilization is pronounced: 25\% more reach with more than 60\% fewer cancelled movements.
Joint training of both levels was additionally shown to degrade performance, establishing that
staged learning is required for this problem.

Three limitations are noted. Training used a single fixed map, leaving room for instance
memorization; evaluation was conducted on instances not used in training, but randomizing instances
during training remains future work. No method completed all 50 deliveries, so the makespan
objective was undefined and delivery count served as a surrogate. Communication was simplified to a
range-based disk model without terrain occlusion.

Future work includes joint fine-tuning of both levels, instance randomization during training, and
measurement of per-step inference latency to quantify the real-time re-decision capability.

\section*{References}

\noindent Agatz, N., Bouman, P., \& Schmidt, M. (2018). Optimization approaches for the traveling
salesman problem with drone. \textit{Transportation Science, 52}(4), 965--981.

\noindent Cesare, K., Skeele, R., Yoo, S.-H., Zhang, Y., \& Hollinger, G. (2015). Multi-UAV
exploration with limited communication and battery. In \textit{Proceedings of the IEEE
International Conference on Robotics and Automation} (pp. 2230--2235).

\noindent Dalal, G., Dvijotham, K., Vecerik, M., Hester, T., Paduraru, C., \& Tassa, Y. (2018).
Safe exploration in continuous action spaces. \textit{arXiv preprint arXiv:1801.08757}.

\noindent Ding, R., Chen, J., Wu, W., Liu, J., Gao, F., \& Shen, X. (2022). Packet routing in
dynamic multi-hop UAV relay network: A multi-agent learning approach. \textit{IEEE Transactions on
Vehicular Technology, 71}(9), 10059--10072.

\noindent Dorling, K., Heinrichs, J., Messier, G. G., \& Magierowski, S. (2017). Vehicle routing
problems for drone delivery. \textit{IEEE Transactions on Systems, Man, and Cybernetics: Systems,
47}(1), 70--85.

\noindent Haarnoja, T., Zhou, A., Abbeel, P., \& Levine, S. (2018). Soft actor-critic: Off-policy
maximum entropy deep reinforcement learning with a stochastic actor. In \textit{Proceedings of the
35th International Conference on Machine Learning} (pp. 1861--1870).

\noindent Huang, S., \& Ontañón, S. (2022). A closer look at invalid action masking in policy
gradient algorithms. \textit{Proceedings of the International FLAIRS Conference, 35}.

\noindent Lowe, R., Wu, Y., Tamar, A., Harb, J., Abbeel, P., \& Mordatch, I. (2017). Multi-agent
actor-critic for mixed cooperative-competitive environments. In \textit{Advances in Neural
Information Processing Systems 30} (pp. 6379--6390).

\noindent Murray, C. C., \& Chu, A. G. (2015). The flying sidekick traveling salesman problem:
Optimization of drone-assisted parcel delivery. \textit{Transportation Research Part C: Emerging
Technologies, 54}, 86--109.

\noindent Murray, C. C., \& Raj, R. (2020). The multiple flying sidekicks traveling salesman
problem: Parcel delivery with multiple drones. \textit{Transportation Research Part C: Emerging
Technologies, 110}, 368--398.

\noindent Ng, A. Y., Harada, D., \& Russell, S. (1999). Policy invariance under reward
transformations: Theory and application to reward shaping. In \textit{Proceedings of the 16th
International Conference on Machine Learning} (pp. 278--287).

\noindent Park, J., Kwon, C., \& Park, J. (2023). Learn to solve the min-max multiple traveling
salesmen problem with reinforcement learning. In \textit{Proceedings of the 22nd International
Conference on Autonomous Agents and Multiagent Systems} (pp. 878--886).

\noindent Yanmaz, E. (2022). Positioning aerial relays to maintain connectivity during drone team
missions. \textit{Ad Hoc Networks, 128}, 102800.

\noindent Zeng, Y., Zhang, R., \& Lim, T. J. (2016). Throughput maximization for UAV-enabled mobile
relaying systems. \textit{IEEE Transactions on Communications, 64}(12), 4983--4996.

\noindent Zhou, Y., Barnes, C., Lu, J., Yang, J., \& Li, H. (2019). On the continuity of rotation
representations in neural networks. In \textit{Proceedings of the IEEE/CVF Conference on Computer
Vision and Pattern Recognition} (pp. 5745--5753).

\end{document}
